\famichapter{Conclusion and Future Scope}{This chapter summarises what was
delivered, what the verification showed, and the planned future work.}{ch:conclusion}

\section{Summary of the Delivered System}

FamiPet delivers a complete, authenticated pet-care and adoption platform. The
backend, built with Node.js, Express, MongoDB, and Mongoose, exposes sixteen
route modules that cover account management with e-mail verification and secure
password recovery; pet management with a unique digital pet identity and QR
code; health records and vaccinations; veterinary appointment booking with
double-booking protection and automatic notification and reminder creation;
care reminders; favourites; community posts with likes and comments; lost-and-found
reports; adoption applications with administrator approval; a per-user
notification hub; and an AI pet-care assistant grounded in the user's real pet
records.

The system enforces the rules of the project brief on the backend: user identity
comes from the signed token, every user-scoped query binds to the authenticated
user, administrative actions are gated by role, data is persisted in MongoDB
and read back through the API, and validation errors are returned as explicit
HTTP responses. The Step 10 verification campaign confirmed persistence,
ownership isolation, role-based 403s, bcrypt password hashing, JWT handling,
CORS behaviour, and a clean browser sweep over all 31 screens.

\section{Observations}

Three implementation choices deserve emphasis because they affect future work.
First, the AI assistant is the only feature that depends on an external paid
service; it falls back to deterministic guidance when the service is
unavailable, so the platform never breaks. Second, notifications and reminders
are produced synchronously by the controllers, which keeps the system simple and
fully consistent at request time, at the cost of no background scheduling or
real-time push. Third, the React client consumes the backend API as the single
source of truth, so the pages always present exactly the persisted state the
backend returns.

\section{Future Scope}

The planned future work is recorded in the roadmap. The items below are planned,
not implemented.

\begin{center}
\begin{longtable}{@{}p{5.2cm}p{9.4cm}@{}}
\caption{Planned future work.}
\label{tab:planned}\\
\toprule
\textbf{Area} & \textbf{Planned work} \\
\midrule
API integration layer & Phase 23: formal typed API client for the React app. \\
Auth and state management & Phase 24: consolidation of authentication and state
management across migrated pages. \\
UI/UX completion & Phase 25: finish the remaining interface surfaces. \\
Visual regression & Phase 26: pixel-level regression for the frontend UI. \\
Functional regression & Phase 27: end-to-end regression across all features. \\
Docker / Nginx & Phase 28: containerization and reverse-proxy deployment. \\
Client consolidation & Phase 29: decommission the legacy browser client once
regression phases are green. \\
Background scheduling & Considered: convert reminders into scheduled jobs and add
push notifications using the unused \texttt{node-cron} and \texttt{socket.io}
dependencies. \\
Provider expansion & Considered: implement the OpenAI adapter described in
\texttt{.env.example} as an alternative AI provider behind the Google default. \\
\bottomrule
\end{longtable}
\end{center}

\section{Closing Statement}

The black book documents the FamiPet project as it exists in the repository:
its requirements, architecture, data model, security, features, implementation,
testing evidence, and development process. Implemented, planned, and considered
items have been clearly distinguished throughout. The system is feature-complete
on the backend, and the React frontend implements the delivered roadmap phases;
the backend remains the source of truth for all
data and authorization.